\documentclass[11pt]{amsart}
\usepackage{amsmath,amssymb,amsthm,mathtools}
\usepackage[T1]{fontenc}
\usepackage{lmodern}
\usepackage{microtype}
\usepackage{booktabs}
\usepackage{graphicx}
\usepackage[hidelinks]{hyperref}

\newtheorem{theorem}{Theorem}[section]
\newtheorem{lemma}[theorem]{Lemma}
\newtheorem{corollary}[theorem]{Corollary}
\newtheorem{proposition}[theorem]{Proposition}
\theoremstyle{remark}
\newtheorem{remark}[theorem]{Remark}

\DeclareMathOperator{\ara}{ara}
\DeclareMathOperator{\adj}{adj}
\DeclareMathOperator{\Min}{Min}

\title{Moh's $P_3$ Curve Is a Global Set-Theoretic Complete Intersection}
\date{September 6, 2026}

\subjclass[2020]{Primary 13C40, 14M10; Secondary 13H05, 13D02, 13P10}
\keywords{set-theoretic complete intersection, arithmetic rank, Moh prime, Hilbert--Burch theorem, symmetric determinantal presentation, explicit equations}

\begin{document}

\begin{abstract}
A classical question asks whether every irreducible affine curve $C'\subseteq\mathbb A^3_{\mathbb C}$ can be written set-theoretically as $C'=V(f,g)$ for two polynomials.  We do not resolve that problem in general.  We prove the global two-equation statement for the singular Moh curve
\[
 t\longmapsto (t^6+t^{31},t^8,t^{10}).
\]
Let $k$ be a field of characteristic zero, let
\[
 \rho:k[x,y,z]\longrightarrow k[t],\qquad
 (x,y,z)\longmapsto(t^6+t^{31},t^8,t^{10}),
\]
and put $\mathfrak p=\ker\rho$.  We construct two explicit degree-$16$ polynomials $H_1,H_2\in\mathbb Z[x,y,z]$ and prove the containment
\[
 \mathfrak p^2\subseteq(H_1,H_2)\subseteq\mathfrak p.
\]
Consequently $\sqrt{(H_1,H_2)}=\mathfrak p$ in the global polynomial ring and $\ara(\mathfrak p)=2$.  In particular the corresponding formal prime in $k[[x,y,z]]$ is a set-theoretic complete intersection, answering the formal-local question recorded by D'Cruz.  The construction starts from the Hilbert--Burch matrix of $\mathfrak p$: a polynomial row operation of constant determinant converts its first three rows into a symmetric $3\times3$ block.  A determinant and an adjugate quadratic form then give $H_1,H_2$, while a symmetric adjugate identity gives the square containment.  The ancillary files verify the parametrization kernel and the polynomial identities used in the two containments by exact arithmetic.
\end{abstract}

\maketitle

\section{Introduction}
The problem that motivates this paper is the following.  Let $C'\subseteq\mathbb A^3_{\mathbb C}$ be an irreducible affine curve and let $I(C')\subseteq\mathbb C[x,y,z]$ be its defining prime ideal.  Since $I(C')$ has height two, one may ask whether there always exist two polynomials $f,g\in\mathbb C[x,y,z]$ such that
\[
 C'=V(f,g),\qquad\text{equivalently}\qquad I(C')=\sqrt{(f,g)}.
\]
In terms of arithmetic rank, this asks whether every such height-two prime has $\ara(I(C'))=2$.  This global affine-space question is open in characteristic zero for arbitrary singular curves; see the surveys of Lyubeznik and D'Cruz \cite{Lyubeznik1989,Lyubeznik1992,DCruz2020}.

We treat the concrete singular curve
\[
 C=\{(t^6+t^{31},t^8,t^{10}):t\in k\}\subseteq\mathbb A^3_k,
\]
with defining prime
\[
 \mathfrak p=\ker\bigl(k[x,y,z]\longrightarrow k[t]\bigr),\qquad
 (x,y,z)\longmapsto(t^6+t^{31},t^8,t^{10}).
\]
Completing at the origin gives the formal prime
\[
 P=\ker\bigl(k[[x,y,z]]\longrightarrow k[[t]]\bigr),
\]
for the same parametrization.  D'Cruz explicitly asked whether this formal prime is a set-theoretic complete intersection \cite[Question 2.5]{DCruz2020}.  Moh's family provides height-two primes that can require arbitrarily many ideal-theoretic generators \cite{Moh1974}.  For the present $P_3$, Gonzalez and Planas-Vilanova proved that four generators are minimal in characteristic zero \cite{GonzalezPlanas2026}.  Thus this example is not an ideal-theoretic complete intersection.

Our main theorem gives the following global containment.

\begin{theorem}\label{thm:main}
Let $k$ be a field of characteristic zero and let
\[
 S=k[x,y,z],\qquad
 \mathfrak p=\ker\bigl(S\longrightarrow k[t]\bigr),
\]
where $(x,y,z)\mapsto(t^6+t^{31},t^8,t^{10})$.  The explicit polynomials $H_1,H_2$ of Appendix~\ref{app:expanded} satisfy
\[
 \mathfrak p^2\subseteq(H_1,H_2)\subseteq\mathfrak p.
\]
Consequently
\[
 \sqrt{(H_1,H_2)}=\mathfrak p
 \qquad\text{and}\qquad
 \ara_S(\mathfrak p)=2.
\]
\end{theorem}

The formal-local statement is a corollary.

\begin{corollary}\label{cor:formal}
Let $R=k[[x,y,z]]$ and let $P$ be the kernel of the same parametrization into $k[[t]]$.  Then
\[
 \sqrt{(H_1,H_2)R}=P.
\]
In particular $\ara_R(P)=2$.
\end{corollary}

\begin{proof}
It remains to justify that the formal kernel is the extended polynomial
kernel.  Put
\[
 B=S/\mathfrak p
   =k[t^6+t^{31},t^8,t^{10}]
 \quad\text{and}\quad
 \mathfrak n=(x,y,z)B.
\]
The ring $k[t]$ is finite over $B$, since $t$ is integral over $B$ by
$t^8=y$.  Moreover $\operatorname{Frac}(B)=k(t)$: in
$\operatorname{Frac}(B)$ one has
\[
 t^2=\frac zy,\qquad
 t^6=\left(\frac zy\right)^3,\qquad
 t^{31}=x-\left(\frac zy\right)^3,
 \qquad
 t=\frac{t^{31}}{(t^2)^{15}}.
\]
Thus the normalization of $B_{\mathfrak n}$ is $k[t]_{(t)}$; the prime
$(t)$ is the unique prime above $\mathfrak n$, since $t^8=y\in\mathfrak n$.
The finite injection
\[
 B_{\mathfrak n}\longrightarrow k[t]_{(t)}
\]
remains injective after $\mathfrak n$-adic completion.  On the right,
$\mathfrak n k[t]_{(t)}=(t^6)$ because $t^6+t^{31}=t^6(1+t^{25})$,
so this completion is $k[[t]]$.  On the left it is
$R/\mathfrak pR$.  The resulting injection is precisely the formal
parametrization, and hence $P=\mathfrak pR$.

The inclusions in Theorem~\ref{thm:main} remain true after extension from
$S$ to $R$.  Taking radicals now gives the claim.
\end{proof}

The proof of Theorem~\ref{thm:main} uses a symmetric presentation of the Hilbert--Burch matrix.  Section~\ref{sec:HB} recalls the four generators and their Hilbert--Burch matrix.  Section~\ref{sec:sym} gives a row operation with constant determinant $216$ that converts the top $3\times3$ block into a symmetric matrix $M$.  If $r$ is the last row, then the maximal minors are generated by
\[
 \Delta=\det M,\qquad v=\adj(M)r^{\mathsf T}.
\]
The two equations are scalar multiples of
\[
 \Delta\qquad\text{and}\qquad \Gamma=r\adj(M)r^{\mathsf T}.
\]
A $3\times3$ symmetric adjugate identity shows that every product $v_iv_j$ lies in $(\Delta,\Gamma)$, hence the square of the prime lies in the two-generated ideal.  No localization, saturation, or primary decomposition is required.

\section{The four generators and the Hilbert--Burch matrix}\label{sec:HB}
Throughout, $S=k[x,y,z]$ with $\operatorname{char}k=0$.  The characteristic-zero generators from \cite{GonzalezPlanas2026} are
\begin{align*}
 f_1={}&x^4-x^2y^4z^3-2xy^6z^2-4xyz-3y^3z^5+3y^3,\\
 f_2={}&x^3y-xy^5z^3-3xz^2-2y^7z^2+2y^2z,\\
 f_3={}&2x^3z-3x^2y^2-2xy^4z^4-y^6z^3+yz^2,\\
 f_4={}&x^2yz-2xy^3-y^5z^4+z^3.
\end{align*}
Thus
\[
 \mathfrak p=(f_1,f_2,f_3,f_4).
\]

Consider the $4\times3$ matrix
\[
\Phi=
\begin{pmatrix}
0&2xy&-2z(2xy^3z^2+y^5z-1)\\
z&-2x^2&y(4x^2yz^3+2xy^3z^2-3)\\
-2y&-3z&-x+6y^2z^4\\
3x&3y&0
\end{pmatrix}.
\]

\begin{lemma}\label{lem:HB}
The signed $3\times3$ minors of $\Phi$, obtained by deleting rows $1,2,3,4$, are respectively
\[
 6f_1,-6f_2,6f_3,-6f_4.
\]
Consequently $I_3(\Phi)=\mathfrak p$, and $\Phi$ is a Hilbert--Burch matrix for $\mathfrak p$.
\end{lemma}

\begin{proof}
The minor identities are direct polynomial calculations.  Since $\mathfrak p$ is a height-two prime and $6\in k^\times$, the ideal of maximal minors is exactly $\mathfrak p$.  The Hilbert--Burch theorem gives the asserted presentation; see \cite[Theorem 20.15]{Eisenbud1995}.
\end{proof}

\section{A constant-determinant symmetrization}\label{sec:sym}
Set
\[
 \alpha=16x^3yz^2-10x^2y^3z+18y^2z^3,
 \qquad
 \beta=16x^2y^2z^2-10xy^4z,
\]
and define
\[
L=
\begin{pmatrix}
0&4&0&0\\
-9&0&0&-\frac83x\\
\alpha&\beta&6&6y^4z^2\\
0&0&0&1
\end{pmatrix}.
\]
A direct calculation gives
\[
 \det L=216.
\]
Hence $L\in\operatorname{GL}_4(S)$.  Multiplying the Hilbert--Burch matrix on the left gives
\[
 L\Phi=
 \begin{pmatrix}
 M\\ r
 \end{pmatrix},
 \qquad r=(3x,3y,0),
\]
where $M$ is the symmetric matrix
\begin{center}
\resizebox{\textwidth}{!}{$\displaystyle
M=
\begin{pmatrix}
4z&-8x^2&4y(4x^2yz^3+2xy^3z^2-3)\\
-8x^2&-26xy&18z(2xy^3z^2+y^5z-1)\\
4y(4x^2yz^3+2xy^3z^2-3)&18z(2xy^3z^2+y^5z-1)&m_{33}
\end{pmatrix},
$}
\end{center}
with
\[
 m_{33}=2\bigl(16x^3yz^3-34x^2y^3z^2-36xy^5z^6+15xy^5z-3x-18y^7z^5+36y^2z^4\bigr).
\]
Because left multiplication by $L$ is an automorphism of $S^4$, the ideal of maximal minors is unchanged:
\[
 I_3\!\left(\begin{matrix}M\\r\end{matrix}\right)=I_3(\Phi)=\mathfrak p.
\]

Put
\[
 \Delta=\det M,\qquad
 v=\adj(M)r^{\mathsf T}=(v_1,v_2,v_3)^{\mathsf T},\qquad
 \Gamma=r\adj(M)r^{\mathsf T}.
\]
Up to signs, the four maximal minors of the $4\times3$ matrix $\left(\begin{smallmatrix}M\\r\end{smallmatrix}\right)$ are $v_1,v_2,v_3,\Delta$.  Hence
\begin{equation}\label{eq:p-vdelta}
 \mathfrak p=(\Delta,v_1,v_2,v_3).
\end{equation}

The following elementary identity is the algebraic core of the proof.

\begin{lemma}[symmetric adjugate identity]\label{lem:adj}
Let $A$ be any commutative ring, let $M\in M_3(A)$ be symmetric, and let $r=(r_1,r_2,r_3)$.  Put
\[
 \Delta=\det M,\qquad C=\adj(M),\qquad v=Cr^{\mathsf T},\qquad \Gamma=rCr^{\mathsf T},
\]
and let
\[
 R_r=
 \begin{pmatrix}
0&-r_3&r_2\\
r_3&0&-r_1\\
-r_2&r_1&0
\end{pmatrix}.
\]
Then
\begin{equation}\label{eq:adj-identity}
 vv^{\mathsf T}-\Gamma C=-\Delta\,R_rMR_r^{\mathsf T}.
\end{equation}
Consequently $v_iv_j\in(\Delta,\Gamma)$ for all $1\le i,j\le3$.
\end{lemma}

\begin{proof}
Equation~\eqref{eq:adj-identity} is the $3\times3$ form of Jacobi's complementary-minor identity.  One may verify it over the universal polynomial ring in the six independent entries of a symmetric $M$ and the three entries of $r$: for $\Delta\ne0$, substitute $C=\Delta M^{-1}$ and apply Jacobi's formula for the $2\times2$ minors of $M^{-1}$; clearing denominators gives \eqref{eq:adj-identity}, hence the polynomial identity holds universally.  Entrywise, the left side therefore differs from a multiple of $\Gamma$ by a multiple of $\Delta$, which gives the final assertion.
\end{proof}

\section{The two global equations}\label{sec:global}
Define
\begin{equation}\label{eq:H-det}
 H_1=-\frac1{16}\Delta,
 \qquad
 H_2=\frac1{36}\Gamma.
\end{equation}
The determinant calculation shows that these are in fact polynomials with integer coefficients; their expanded forms are recorded in Appendix~\ref{app:expanded}.  Both have degree $16$.

\begin{proposition}\label{prop:containment}
The polynomials in \eqref{eq:H-det} satisfy
\[
 \mathfrak p^2\subseteq(H_1,H_2)\subseteq\mathfrak p.
\]
\end{proposition}

\begin{proof}
By \eqref{eq:p-vdelta}, $\Delta\in\mathfrak p$ and each $v_i\in\mathfrak p$.  Since
\[
 \Gamma=r_1v_1+r_2v_2+r_3v_3,
\]
we also have $\Gamma\in\mathfrak p$.  Thus $(H_1,H_2)=(\Delta,\Gamma)\subseteq\mathfrak p$, because $16$ and $36$ are units.

For the reverse containment, \eqref{eq:p-vdelta} shows that $\mathfrak p^2$ is generated by products among $\Delta,v_1,v_2,v_3$.  Every product containing $\Delta$ lies in $(\Delta,\Gamma)$, and Lemma~\ref{lem:adj} gives
\[
 v_iv_j\in(\Delta,\Gamma)
 \qquad(1\le i,j\le3).
\]
Hence $\mathfrak p^2\subseteq(\Delta,\Gamma)=(H_1,H_2)$.
\end{proof}

\begin{proof}[Proof of Theorem~\ref{thm:main}]
Proposition~\ref{prop:containment} gives
\[
 \mathfrak p^2\subseteq(H_1,H_2)\subseteq\mathfrak p.
\]
Since $\mathfrak p$ is prime,
\[
 \mathfrak p=\sqrt{\mathfrak p^2}
 \subseteq\sqrt{(H_1,H_2)}
 \subseteq\sqrt{\mathfrak p}=\mathfrak p.
\]
Thus $\sqrt{(H_1,H_2)}=\mathfrak p$.  Since $\operatorname{ht}\mathfrak p=2$, the arithmetic rank is exactly two.
\end{proof}

\begin{remark}\label{rem:charp}
The containment identities in Section~\ref{sec:verification} use only integer coefficients and scalar factors divisible by $2$ and $3$.  Combined with the generator theorem of \cite{GonzalezPlanas2026}, the same global conclusion therefore persists in characteristic $p\ge5$.  We state the main theorem in characteristic zero because that is the case of the classical open question considered here.
\end{remark}

\section{Comparison with the local construction}
The global equations $H_1,H_2$ are different from the degree-$27$ and degree-$22$ pair obtained from the conormal/canonical-module construction in the earlier local argument.  That local pair correctly defines the formal germ at the origin but possesses the additional global component
\[
 (x,4z^5-1).
\]
For the present pair, the global containment
\[
 \mathfrak p^2\subseteq(H_1,H_2)\subseteq\mathfrak p
\]
precludes additional minimal primes.  No localization or saturation is required.

The proof uses a structural feature of the Hilbert--Burch matrix.  The constant-determinant row operation $L$ turns three of its rows into a symmetric matrix.  For a $4\times3$ matrix of the form $\left(\begin{smallmatrix}M\\r\end{smallmatrix}\right)$ with $M$ symmetric, the determinant $\det M$ and the adjugate quadratic form $r\adj(M)r^{\mathsf T}$ automatically generate an ideal containing the square of the maximal-minor ideal.  In the Moh example the symmetrization occurs globally over the polynomial ring because $\det L=216$ is constant.

\section{Computational verification}\label{sec:verification}
The proof does not require a direct radical computation.  The explicit polynomial identities can also be checked by exact computer algebra.

Over $\mathbb Q$, the ancillary verification package performs the following checks.
\begin{enumerate}
\renewcommand{\labelenumi}{(\roman{enumi})}
\item It recomputes the kernel of
\[
 \mathbb Q[x,y,z]\longrightarrow\mathbb Q[t],\qquad
 (x,y,z)\longmapsto(t^6+t^{31},t^8,t^{10}),
\]
by elimination and verifies that it equals $(f_1,f_2,f_3,f_4)$.
\item It verifies $\det L=216$ and the exact matrix identity
\[
 L\Phi=\begin{pmatrix}M\\r\end{pmatrix},\qquad M=M^{\mathsf T}.
\]
\item It verifies
\[
 \det M=-16H_1,
 \qquad
 r\adj(M)r^{\mathsf T}=36H_2.
\]
\item It verifies the short membership identities
\begin{align*}
H_1={}&-24x f_1+\bigl(64x^4yz^2-40x^3y^3z+72xy^2z^3-81y^4z^2\bigr)f_3+81f_4,\\
H_2={}&-12y f_1+27x f_2+\bigl(-40x^3y^2z^2+25x^2y^4z+36y^3z^3\bigr)f_3,
\end{align*}
which directly prove $(H_1,H_2)\subseteq\mathfrak p$.
\item For each $1\le i\le j\le4$, it supplies integer polynomials $A_{ij},B_{ij}$ and a positive integer $d_{ij}$ such that
\[
 d_{ij}f_if_j=A_{ij}H_1+B_{ij}H_2.
\]
The ten scalar factors are
\[
\begin{array}{c|cccccccccc}
(i,j)&11&12&13&14&22&23&24&33&34&44\\ \hline
d_{ij}&5832&2916&972&8748&1458&486&4374&162&1458&13122.
\end{array}
\]
As a compact example,
\[
 162f_3^2=-y(5x^3-2yz)H_1-x(8x^3+13yz)H_2.
\]
These ten identities prove $\mathfrak p^2\subseteq(H_1,H_2)$ over characteristic zero.
\end{enumerate}

The supplied checker uses sparse rational polynomial arithmetic and verifies the displayed identities independently of the script that produced them.  A Macaulay2 replay file is also included so that the kernel computation and the ideal containments can be reproduced in a standard commutative-algebra system \cite{Macaulay2}.  No floating-point computation or numerical root finding is used.

\appendix
\section{Expanded global equations}\label{app:expanded}
We record the two global equations in expanded form:
\begin{align*}
H_1={}&128x^7yz^3-272x^6y^3z^2-128x^5y^5z^6+120x^5y^5z-24x^5\\
&+16x^4y^7z^5+208x^4y^2z^4+40x^3y^9z^4-394x^3y^4z^3\\
&-144x^2y^6z^7+291x^2y^6z^2+177x^2yz+90xy^8z^6\\
&+144xy^3z^5-234xy^3+81y^{10}z^5-162y^5z^4+81z^3,
\end{align*}
\begin{align*}
H_2={}&-80x^6y^2z^3+170x^5y^4z^2+80x^4y^6z^6-75x^4y^6z+15x^4y\\
&-10x^3y^8z^5+32x^3y^3z^4-25x^2y^{10}z^4-98x^2y^5z^3-81x^2z^2\\
&-72xy^7z^7-30xy^7z^2+102xy^2z-36y^9z^6+72y^4z^5-36y^4.
\end{align*}

\section{Verification files}
The ancillary package contains the expanded equations, the ten square-containment identities, machine-readable coefficient data, a standalone exact-arithmetic checker, and a Macaulay2 replay script.  The radical equality follows directly from the two ideal containments and does not require a primary decomposition.

\clearpage
\thispagestyle{plain}
\section*{Acknowledgments}
This work grew out of a study of Moh's $P_3$ prime and benefited from exact computational experiments in Macaulay2 and formal verification in Lean~4.  The author is grateful for the mathematical discussions and feedback that helped shape the presentation.

\section*{Funding}
No external funding is reported for this work.

\section*{Declaration of generative AI and AI-assisted technologies in the manuscript preparation process}
During the mathematical development and preparation of this manuscript, the author used OpenAI's ChatGPT and Codex and Harmonic's Aristotle for symbolic experimentation, proof organization, language revision, and assistance in constructing and checking a Lean~4 formalization.  Macaulay2 was used for independent exact-arithmetic computational checks.  The author remains responsible for the mathematical content and final presentation.

\clearpage
\begin{thebibliography}{99}

\bibitem{Boratyński1978}
M.~Boraty\'nski,
\emph{A note on set-theoretic complete intersection ideals},
J. Algebra \textbf{54} (1978), no.~1, 1--5.

\bibitem{Bresinsky1979}
H.~Bresinsky,
\emph{Monomial space curves in $\mathbb A^3$ as set-theoretic complete intersections},
Proc. Amer. Math. Soc. \textbf{75} (1979), no.~1, 23--24.

\bibitem{CowsikNori1978}
R.~C. Cowsik and M.~V. Nori,
\emph{Affine curves in characteristic $p$ are set-theoretic complete intersections},
Invent. Math. \textbf{45} (1978), no.~2, 111--114.

\bibitem{DCruz2020}
C.~D'Cruz,
\emph{Symbolic powers, set-theoretic complete intersection and certain invariants},
The Special Issue of the Proceedings of the Telangana Academy of Sciences \textbf{1} (2020), no.~1, 75--84; arXiv:2003.09101.

\bibitem{Eisenbud1995}
D.~Eisenbud,
\emph{Commutative Algebra with a View Toward Algebraic Geometry},
Graduate Texts in Mathematics, vol.~150, Springer, New York, 1995.

\bibitem{GonzalezPlanas2026}
L.~Gonz\'alez and F.~Planas-Vilanova,
\emph{Prime ideals of Moh and the characteristic of the field},
J. Algebra Appl., published online 17 March 2026, article 2750170; also arXiv:2407.21692.

\bibitem{Lyubeznik1989}
G.~Lyubeznik,
\emph{A survey of problems and results on the number of defining equations},
in \emph{Commutative Algebra} (Berkeley, CA, 1987), Math. Sci. Res. Inst. Publ., vol.~15, Springer, New York, 1989, pp.~375--390.

\bibitem{Lyubeznik1992}
G.~Lyubeznik,
\emph{The number of defining equations of affine algebraic sets},
Amer. J. Math. \textbf{114} (1992), no.~2, 413--463.

\bibitem{Macaulay2}
D.~R. Grayson and M.~E. Stillman,
\emph{Macaulay2, a software system for research in algebraic geometry},
\url{https://macaulay2.com/}.

\bibitem{Moh1974}
T.~T. Moh,
\emph{On the unboundedness of generators of prime ideals in power series rings of three variables},
J. Math. Soc. Japan \textbf{26} (1974), no.~4, 722--734.

\bibitem{Moh1982}
T.~T. Moh,
\emph{A result on the set-theoretic complete intersection problem},
Proc. Amer. Math. Soc. \textbf{86} (1982), no.~1, 19--20.

\bibitem{Moh1985}
T.~T. Moh,
\emph{Set-theoretic complete intersections},
Proc. Amer. Math. Soc. \textbf{94} (1985), no.~2, 217--220.

\end{thebibliography}

\end{document}
